\section{Introduction}

A large organisation typically buys from tens of thousands of distinct supplier
records. Many of those records are not distinct suppliers. They are divisions, acquired
subsidiaries, regional entities, and renamed business units of a much smaller number of
ultimate parents. Until those records are grouped into corporate families, a buyer
cannot answer the question that motivates the exercise: how much do we actually spend
with this company, and what leverage does that give us? The same grouping underpins
credit exposure aggregation, beneficial-ownership and sanctions screening, and supplier
risk consolidation, and the same structure appears on the revenue side as account
hierarchy in customer systems.

This task is routinely treated as an instance of entity resolution, the problem of
deciding whether two records refer to the same real-world
entity~\citep{papadakis2021fourgenerations,barlaug2021survey}. The framing is convenient
because entity resolution has mature benchmarks and a decade of strong neural
methods~\citep{mudgal2018deepmatcher,li2020ditto,peeters2025entitymatchingllm}. It is also wrong
in a way that matters. Entity resolution asks whether two records are the \emph{same}
thing. Corporate-family resolution asks whether two records that are emphatically
\emph{different} things, separate legal entities with separate registrations and often
separate countries, belong to the same owner. The two questions have different
evidence. For duplicate detection, the evidence is in the records: names, addresses,
and identifiers of the same entity tend to look alike. For family linkage, the evidence
is frequently in neither record. That \textsc{Heico Corp} owns \textsc{Blue Aerospace
LLC} is a fact recorded in a corporate filing. It is not a fact about the two strings.

If that is true, a benchmark that mixes both kinds of pair and reports a single number
will systematically mislead. We introduce \textsc{CorpFam}, a public benchmark of
\BenchPairs candidate pairs over \BenchFamilies corporate families built from
\ArchiveRows federal award records, and stratify every pair by \emph{name visibility}:
whether the two names are identical after normalisation, share a distinctive token, or
share none at all.

The stratification changes the conclusion, and it does so twice.

First, at the matching stage. We report per-stratum \emph{recall} rather than $F_1$, for
a reason that turns out to matter a great deal: the strata have positive rates ranging
from \BaseVisibleRate to \BaseIdenticalRate, so a classifier that answers yes to every
pair scores \BaseIdenticalF on identical pairs and \BaseVisibleF on visible ones without
doing anything at all. Per-stratum $F_1$ therefore measures the base rate more than the
method, whereas recall depends only on a stratum's own positives. On that footing the
result is unambiguous: the best matcher recovers \BestRecallIdentical of identical
pairs, \BestRecallVisible of visible pairs, and \BestRecallInvisible of invisible ones
(95\% CI \BestRecallInvisibleCI). No method exceeds \AnyInvisibleRecallMax.

Second, and this is the finding we think matters most, the failure does not begin
at matching. Blocking decides which pairs a matcher ever sees. Part of this is true by
definition and we say so plainly: the invisible stratum is defined by an empty
intersection of distinctive tokens, and token blocking keys on exactly that function, so
its \BlkTokenInvisible recall there is an identity rather than a discovery. The
empirical question is what happens once the key is moved away from that function. We
therefore also evaluate character q-grams, sorted-neighbourhood adjacency, phonetic
keys, attribute keys that ignore the name entirely, and semantic nearest neighbours in
an embedding space. None exceeds three percent on the invisible stratum: embedding
nearest-neighbour blocking reaches \BlkAnnInvisible and attribute blocking
\BlkAttrInvisible. Unioning every scheme recovers \BlkUnionInvisible, leaving
\BlkInvisibleLost of these links outside the candidate set before any matcher runs. A
better model cannot classify a pair it never sees.

A reviewer might reasonably suspect that a stratum defined by dissimilar names and then
shown to defeat name-based methods is circular, or that its links are registry errors.
We address both. Against SEC Exhibit~21 subsidiary schedules, filed under securities
law by the parent's own counsel and sharing no provenance with procurement registration,
\EdgarInvisibleConfirmPct of invisible links are corroborated, against
\EdgarRandomParentPct when the identical procedure is run against permuted parents, an
excess over chance of \EdgarExcessOverChance points. And a supervised model given
address, geography and phone features does not rescue the stratum either, for a
structural reason: only \ParentsTransactingPct of parents ever appear as an award
recipient, so for the rest there is no address to compare against.

\paragraph{Contributions.}
\begin{enumerate}
\item \textbf{A public benchmark for corporate-family resolution}: \BenchPairs pairs
over \BenchFamilies families, with ground truth that is self-reported to a government
registry rather than hand-annotated or model-derived (Section~\ref{sec:benchmark}).

\item \textbf{A name-visibility stratification} that we argue is the correct primary
axis of evaluation for this task, and without which aggregate metrics are not
interpretable (Section~\ref{sec:strata}).

\item \textbf{Evidence that the standard pipeline fails at blocking, not matching.}
\BlkInvisibleLost of invisible links are never proposed as candidates, including by
blockers that are not functions of token overlap: phonetic, attribute-keyed, and
embedding nearest-neighbour (Section~\ref{sec:blocking}).

\item \textbf{Independent corroboration of the hard stratum} against SEC Exhibit~21,
with a permutation control establishing that corroboration exceeds chance by
\EdgarExcessOverChance points and that these are genuine corporate relationships rather
than label noise (Section~\ref{sec:edgar}).

\item \textbf{Baselines on a single protocol}: \NumMethods pairwise matchers spanning
string similarity and sentence embeddings, supervised multi-attribute models over
address, geography and phone features, and a family-level clustering task
(Sections~\ref{sec:results}--\ref{sec:clustering}).
\end{enumerate}

\section{Related Work}

\paragraph{Entity resolution benchmarks.} The standard suites established the task and
drove rapid progress~\citep{li2020ditto,zeakis2023pretrained}: Magellan and
DeepMatcher~\citep{konda2016magellan,mudgal2018deepmatcher}, WDC
Products~\citep{peeters2024wdcproducts}, Alaska~\citep{crescenzi2021alaska}, and
Machamp~\citep{wang2021machamp}. Several contain a company or organisation
split, and recent work extends evaluation to large language
models~\citep{peeters2023chatgpt,zhang2024jellyfish,peeters2025entitymatchingllm} and to
unified cross-task models~\citep{tu2023unicorn}. Critiques have noted that easy
negatives inflate reported
performance~\citep{papadakis2023benchmarking,neuhof2024openbenchmark}, and that skewed positive-class
ratios make reported scores incomparable across benchmarks~\citep{li2022bridging}, a
hazard we take seriously enough that Section~\ref{sec:whyrecall} is devoted to it. Our
concern is adjacent but distinct: not that the negatives are too easy, but that the
\emph{positives} are heterogeneous in a way a single metric cannot express, because some
are solvable from the records and some are not solvable from the records at all.
MaDI-Bench~\citep{steiner2026madibench} moves in a similar direction by systematising
corner-case and blocking difficulty across an integration pipeline; the difference is
that its hard positives are synthesised as surface variants, whereas ours occur naturally
and are certified by an external legal filing.

\paragraph{Company graphs.} CompanyKG~\citep{cao2023companykg} is the closest public
resource, a large heterogeneous graph of company relationships with an edge-prediction
task that includes acquisition edges. It is complementary rather than overlapping: its
nodes are already resolved company entities. The linkage problem we study, deciding
whether two \emph{registration records} belong to one family given only what a
procurement system actually holds, is assumed away before its task begins.

\paragraph{Blocking and candidate generation.} Blocking is the step that makes entity
resolution tractable, and it has a large literature of its own: token and q-gram
blocking, sorted neighbourhood, canopies, meta-blocking and learned
blockers~\citep{papadakis2021fourgenerations,thirumuruganathan2021blocking}. Evaluation is
conventionally in terms of pair completeness and reduction ratio. What we add is the
observation that on a task where the relationship is not encoded in the compared
strings, every scheme in this family shares a single blind spot, and that measuring
blocking per difficulty stratum exposes a ceiling that aggregate pair completeness
hides.

\paragraph{Company name matching and harmonisation.} A long line of work matches
company names across sources, in patent
data~\citep{magerman2006patentee,thoma2010harmonizing,peeters2010harmonizing}, in financial
identifiers~\citep{chan2013lei,chan2019lei,arimond2023elf}, and with learned name
representations~\citep{ziv2022companyname2vec}. This literature is largely concerned
with recognising the \emph{same} company under spelling variation.
Corporate-family resolution begins where that ends.

\paragraph{Corporate hierarchies and ownership.} Ownership networks have been studied in
economics and network
science~\citep{vitali2011network,mizuno2020corporatecontrol,garciabernardo2017offshore}, in bank
holding structures~\citep{flood2020bhc}, and in knowledge graphs of European business
registers~\citep{roman2022eubusinessgraph,soylu2022theybuyforyou}. Hierarchical linkage
has been treated as link prediction over enterprise
data~\citep{ebeid2021hierarchical,ganesan2020linkprediction,ganesan2024xlp}. These works build
or analyse hierarchies; to our knowledge none releases a public, reproducible benchmark
with held-out splits, and none stratifies by whether the relationship is recoverable
from the records.

\paragraph{Procurement and master data.} The practical stakes are documented in work on
supplier master data
quality~\citep{athanasiadou2019suppliermdm,ibrahim2021masterdataquality,gantman2021vendormaster},
duplicate invoice detection~\citep{ho2009duplicateinvoices,liao2026procurementduplicate},
and purchasing consolidation
savings~\citep{smart2007purchasingsynergy,carril2020consolidation}.

\section{The \textsc{CorpFam} Benchmark}
\label{sec:benchmark}

\subsection{Source and ground truth}

We build on United States federal contract award data, published in bulk and in
full~\citep{usaspending_api}. Each award record identifies the recipient by a Unique Entity
Identifier (UEI) and separately records the recipient's \emph{ultimate parent} UEI and
name. That field is not an inference by us or by the publisher: it is self-reported by
the supplier during federal registration, where misstatement carries legal consequence,
and it is exactly the relation we want to predict. This gives ground truth independent
of any matching system, which distinguishes it from benchmarks whose labels come from
human annotation of the same strings a model will see.

We processed the complete FY2025 contract archive: \ArchiveRows award rows covering
\ArchiveEntities distinct entities.

\subsection{The self-parent trap}
\label{sec:selfparent}

The parent field cannot be used naively. Of the entities reporting a parent at all,
\ArchiveSelfParentPct report \emph{themselves}: the parent UEI equals the entity's own
UEI. This is correct registration behaviour for a company with no corporate parent, but
it means a pipeline treating every populated parent field as a family link would build a
dataset in which the overwhelming majority of ``links'' join an entity to itself. Such
pairs are trivially positive under any string method and would push every reported score
toward the ceiling.

We confirmed the pattern independently in the per-entity API, where \AuditSelfParentPct
of child-level records are self-parenting, agreeing by both UEI equality and internal
identifier stem. Removing self-parents leaves \ArchiveGenuineLinks genuine parent--child
links, the population the benchmark is built from.

\subsection{Name visibility as the primary stratum}
\label{sec:strata}

A second and subtler triviality remains. Many genuine links connect a child to a parent
whose name, after normalisation, is identical: distinct legal registrations under a
shared trading name. These links are genuine and belong in the data, but they are
solvable by string equality and say nothing about a method's grasp of corporate
structure. Of the \ArchiveGenuineLinks genuine links, \NvIdenticalRaw are name-identical,
leaving \NvInformative informative links, of which \NvInvisiblePctInformative share no
distinctive token with their parent.

We therefore assign every pair to one of three strata:

\begin{description}
\item[Identical] the names are equal after case-folding, punctuation removal and
whitespace normalisation;
\item[Visible] the names differ but share at least one distinctive token, where
distinctive excludes legal forms and generic corporate vocabulary (\textsc{inc},
\textsc{holdings}, \textsc{international}, \textsc{services}, and similar);
\item[Invisible] the names share no distinctive token at all.
\end{description}

Table~\ref{tab:composition} gives the composition: \BenchPosIdenticalPct identical,
\BenchPosVisiblePct visible, \BenchPosInvisiblePct invisible. That first figure is why a
single aggregate is not interpretable here. More than half the positives are solvable
by exact match, so a method that only normalises strings starts from a high floor.

\begin{table}[t]
\centering
\small
\begin{tabular}{lrr}
\toprule
Stratum & Positive pairs & \% \\
\midrule
Identical & 7,541 & 55.0 \\
Visible & 2,953 & 21.5 \\
Invisible & 3,222 & 23.5 \\
\midrule
All positives & 13,716 & 100.0 \\
\bottomrule
\end{tabular}
\caption{Composition of \textsc{CorpFam} positive pairs by name visibility.
\emph{Identical} pairs share a normalised name, \emph{visible} pairs share at
least one distinctive token, and \emph{invisible} pairs share none. Negatives are
generated at three per positive, giving 54,864 pairs in total.}
\label{tab:composition}
\end{table}


The invisible stratum is not a curiosity. It contains the cases that motivate the task
commercially: an acquired subsidiary trading under its original name, or a holding
company whose operating units carry unrelated brands. It is also where spend
concentrates, since acquisitive parents are large.

\subsection{Why we report recall, not per-stratum $F_1$}
\label{sec:whyrecall}

Stratifying creates a measurement hazard that is easy to walk into, and we walked into
it before catching it. The strata do not have comparable positive rates: on the test
split they are \BaseIdenticalRate, \BaseVisibleRate and \BaseInvisibleRate respectively.
A classifier that answers yes to every pair therefore scores \BaseIdenticalF on the
identical stratum and \BaseVisibleF on the visible one, purely as a function of
composition. Comparing $F_1$ across strata under those conditions largely compares base
rates, and a large fraction of any apparent gap is an artifact of the denominators.
Skewed positive-class ratios have been shown to make entity-matching results
incomparable in exactly this way~\citep{li2022bridging}.

Worse, the negative grades are strongly associated with stratum by construction: a
hard-string negative shares a distinctive token and is therefore \emph{necessarily}
visible, while a randomly drawn negative almost never shares one and is therefore almost
always invisible. The hard stratum thus receives the easiest negatives, which biases
per-stratum precision in the opposite direction to the effect we are measuring.

We therefore report per-stratum \textbf{recall}, which depends only on the positives of
that stratum and is invariant to both problems, and we report precision and average
precision globally, where the negative pool is shared. Every table carries the
majority-class baseline so no cell can be read without its floor. This is not a
presentational choice: two intermediate results in this project reached $88$ and $59.9$
on the invisible stratum and both were the base rate rather than a finding.

\subsection{Splits, negatives, and leakage}
\label{sec:splits}

\paragraph{Splits are assigned by family, never by pair.} If members of one family appear
in both training and test, a model can memorise the family during training and be scored
on recall of that memory. We assign each family wholly to train, validation or test in a
$60/20/20$ ratio, giving \BenchPosTrain, \BenchPosVal and \BenchPosTest positive pairs,
and the build asserts that no family spans two splits, failing loudly rather than
reconciling silently.

\paragraph{Negatives match the role structure of positives.} Each negative pairs a child
with the parent of a \emph{different} family. This is not cosmetic, and we report it
because the alternative is a trap we initially fell into. An earlier build drew negatives
as child--child pairs. Because parent entities are frequently holding companies that never
transact directly, they carry no address in the award data, so ``both sides have an
address'' separated the classes almost perfectly. A supervised model exploiting that
shortcut reached $88$ $F_1$ on the invisible stratum while learning nothing about
ownership. It had learned to detect entity \emph{role}. Matching the role composition
of positives and negatives removes the shortcut. Any benchmark of hierarchical
relationships is exposed to this failure, since the two sides of a hierarchical edge are
by definition different kinds of object.

\paragraph{Negatives come in three grades}, three per positive, drawn within each split:
\emph{random} (a child and an unrelated family's parent), \emph{hard string} (sharing a
distinctive token with an unrelated parent), and \emph{same block} (sharing a blocking
key, i.e.\ what a deployed pipeline actually adjudicates). Reporting against random
negatives alone is the standard way to overstate
precision~\citep{papadakis2023benchmarking}.

\section{Ground-Truth Quality}
\label{sec:quality}

A benchmark is only as trustworthy as its labels, and we think it more useful to
quantify their limits than to assert their correctness.

\subsection{The registry disagrees with itself}

The same registration is exposed through two independent channels: a per-entity API and
the bulk award archive. On the \ReconBoth entities where both report a distinct parent,
they agree on the exact parent UEI in \ReconAgreeUei of cases, on the normalised parent
name in \ReconAgreeName (Wilson 95\% CI \ReconAgreeNameCI), and on any shared distinctive
token in \ReconAgreeToken. In \ReconDisagreePct of cases the two sources name a
substantively different parent (Table~\ref{tab:reconciliation}). A further \ReconOnlyOne
entities have a parent in one source and none in the other.

\begin{table}[t]
\centering
\small
\begin{tabular}{lr}
\toprule
Agreement criterion & \% \\
\midrule
Exact parent UEI & 64.9 \\
Normalised parent name & 73.8 \\
Any shared distinctive token & 81.5 \\
\midrule
Substantively different parent & 18.5 \\
\bottomrule
\end{tabular}
\caption{Agreement between the two independent views of the same ground truth --
the per-entity API and the bulk award archive -- on the 248
entities where both report a distinct parent. Even the registry disagrees with
itself on roughly one link in five, which bounds the accuracy any method can be
credited with.}
\label{tab:reconciliation}
\end{table}


Two further instabilities are worth stating plainly. Within the archive alone,
\ArchiveConflicting entities are assigned more than one distinct parent across their own
award rows, which we resolve by obligated-value plurality. And in the per-entity API,
\AuditMultiParent of \AuditChildRecords child records carry more than one parent in the
returned history, up to \AuditMaxParents on a single record, because the field
accumulates rather than replaces as ownership changes. We take the most recent, but the
ambiguity is itself a finding about registry-derived ground truth.

Inspection of the disagreements shows three causes: genuine ownership change between
snapshots, registration identity lagging a completed acquisition, and outright error. We
do not arbitrate. The practical consequence is a ceiling: a method credited with much
more than roughly \ReconAgreeName agreement against this ground truth is likely fitting
registry noise.

\subsection{Independent corroboration from SEC filings}
\label{sec:edgar}

The two views above share a provenance. If a supplier misdeclared its parent at
registration, both inherit the error identically, so their agreement bounds internal
consistency without establishing correctness. This matters most for the invisible
stratum, where a sceptic can reasonably propose that links with no name relationship are
simply data-entry errors, which would make our headline finding an artifact of dirty
labels rather than a property of the task.

Exhibit~21 settles the question from outside. It is the ``Subsidiaries of the
Registrant'' schedule filed with a 10-K, prepared by the parent's own counsel under
securities law, and it shares no provenance with a procurement registration. We retrieved
Exhibit~21 for the largest families in the benchmark, obtaining \EdgarFilings usable
filings, and checked \EdgarLinks declared links against them.

\begin{table}[t]
\centering
\small
\begin{tabular}{lrrrr}
\toprule
Stratum & Checked & Confirmed & \% & 95\% CI \\
\midrule
Identical & 157 & 107 & 68.2 & [60.5, 74.9] \\
Visible & 544 & 391 & 71.9 & [68.0, 75.5] \\
Invisible & 692 & 336 & 48.5 & [44.9, 52.3] \\
\midrule
All & 1,393 & 834 & 59.9 &
[57.3, 62.4] \\
\midrule
\multicolumn{5}{l}{\emph{Invisible stratum against controls}} \\
\quad wrong parent, 20 permutations & -- & -- &
0.33 & -- \\
\quad nearest-size wrong parent & -- & -- &
3.90 & -- \\
\bottomrule
\end{tabular}
\caption{Independent corroboration against SEC Exhibit~21 subsidiary schedules over
132 filings. Exhibit~21 is filed under securities law by the
parent's own counsel and shares no provenance with a supplier's procurement
registration. The control rows re-run the identical matcher against the wrong parent's
filing: chance corroboration is near zero, so the invisible-stratum rate is not an
artifact of long filings or a permissive matcher, and the stratum consists of genuine
corporate relationships rather than registry noise. Non-confirmation is weak evidence
in the other direction, since Exhibit~21 omits immaterial subsidiaries and its
formatting is unregulated.}
\label{tab:edgar}
\end{table}


Overall \EdgarConfirmPct of links are corroborated. On the invisible stratum
(\EdgarInvisibleChecked links) \EdgarInvisibleConfirmPct are confirmed, with 95\%
interval \EdgarInvisibleCI. \textsc{Aerojet Rocketdyne} under \textsc{L3Harris
Technologies} is typical: a real acquisition, entirely invisible in the names, confirmed
by filing.

\paragraph{A rate is meaningless without a chance floor.} Exhibit~21 lists are long and
our matcher is fuzzy, so some corroboration would occur by accident. We therefore re-ran
the identical procedure with families matched to the \emph{wrong} parent's filing. Over
\EdgarPermutations permutations, invisible-stratum corroboration against a randomly
assigned parent is \EdgarRandomParentPct. A harder control assigns each family the
filing whose parsed length is closest to its true parent's, removing any advantage from
filing size, and gives \EdgarNearestSizePct. The observed rate therefore exceeds
chance by \EdgarExcessOverChance points, and the corroboration is not an artifact of
long filings or a permissive matcher.

Two caveats belong here rather than in a footnote. Corroboration correlates with how much
text our parser recovered from a filing (Spearman $\rho=\EdgarYieldRho$), which means the
reported rate is attenuated by parser recall: \EdgarInvisibleConfirmPct is a lower bound
on true corroboration, not an inflated figure. And because we resolve companies to SEC
identifiers by name, the filings we obtain are selected on the parent's name being
matchable against the registrant list, so coverage is not a random sample of families.
Non-confirmation is in any case weak evidence in the other direction: Exhibit~21 omits
immaterial subsidiaries and its formatting is unregulated.

\subsection{Adjudicating family artifacts by hand}

Some declared families are not families. The archive asserts, for instance, that an
Australian subsidiary is the ultimate parent of \textsc{Raytheon Company}, and that a
moving-services company is the parent of one hundred and three waste-management entities.

Our first instinct was an automatic rule: flag a family when its children share a dominant
token the declared parent lacks and the parent's own name does not reach its children.
That rule catches the cases above. It also flagged \textsc{Arctic Slope Regional
Corporation}, whose subsidiaries are entirely genuine. This is not a tuning failure. A
rule keyed on name dissimilarity cannot distinguish a mislabelled family from a real
conglomerate, because a real conglomerate looks exactly like a mislabelled family from the
outside. Applied aggressively it would remove \textsc{Heico}, \textsc{TransDigm},
\textsc{Berkshire Hathaway} and \textsc{Republic Services}, precisely the
invisible-stratum families that make the benchmark worth building.

We therefore adjudicated large families by hand against their child lists. We inspected
\BenchInspected families, excluded \BenchExcludedParents and retained \BenchRetained,
and publish the full log, retentions as well as rejections, each with a written
reason. Publishing only the rejections would make a careful audit indistinguishable from
a purge of inconvenient cases, and the retentions are the more informative half: they
record where the automatic rule was wrong. Exclusions are keyed on entity identifier
rather than on name string, which in a paper about name instability is the only
defensible choice. This removes \BenchDroppedArtifact links, plus \BenchDroppedSovereign
under sovereign catch-all entities, and the complete sovereign list of
\BenchSovereignAll parents ships with the manifest rather than a truncated sample of it.

\paragraph{The exclusions are not neutral, and we report their effect.} Of the links
removed by hand adjudication, \BenchDropInvisiblePct fall in the invisible stratum,
the very stratum the paper's argument rests on. That is expected, because a mislabelled
family looks exactly like a name-invisible one, but it means a sceptic is entitled to ask
whether the headline survives the filter being removed. We therefore rebuild the entire
benchmark with no exclusions at all, \SensUnfilteredLinks links of which
\SensUnfilteredInvisible are invisible, and re-run the matching experiments against
it. The best invisible-stratum recall on the unfiltered build is
\SensUnfilteredInvisibleRecall, against \SensFilteredInvisibleRecall on the filtered one.

The conclusion therefore holds \emph{a fortiori}: removing the exclusions makes the task
harder, not easier. Whatever the adjudication did, it did not manufacture the finding by
discarding difficult cases. If anything it discarded cases that were unsolvable for
the wrong reason, and the benchmark is marginally more tractable with them gone. Both
builds ship. Smaller families are retained without automatic filtering in either, and we
report this as a bounded source of label noise rather than one we claim to have
eliminated.

\subsection{Provenance and redistribution}
\label{sec:provenance}

Because the benchmark publishes company names in bulk, we checked whether we may
redistribute them rather than assuming it. The data is public-domain United States
Government work~\citep{usaspending_api}. However, SAM.gov's terms grant only a limited
licence over ``D\&B Open Data'', which explicitly includes legal business name, and forbid
disseminating it in bulk; those terms scope the restriction to records associated with
base awards dated before \DnbCutoff.

The definitive test named in those terms is an Entity Validation Service source field,
which does not exist in the award data. We therefore bounded the exposure. Scanning all
\DnbLinksScanned award rows and taking the earliest performance-period start across the
transactions supporting each link, \DnbPctPre of positive pairs have at least one
supporting award beginning before the cutoff. This is deliberately an upper bound: a task
order placed in FY2025 against a long-running vehicle inherits that vehicle's original
start date. Under the opposite reading, that every transaction here is an FY2025 action,
the exposure is zero. We release the full benchmark and, alongside it, a conservative
post-cutoff subset of \DnbConservativePos positive pairs.

\section{Blocking: The Failure Before Matching}
\label{sec:blocking}

Every result reported so far in the literature, and every result in the next section,
measures a matcher \emph{given} a pair. Deployed entity resolution never sees all pairs.
A blocking step first proposes candidates, and anything it fails to propose is
unrecoverable regardless of the matcher. On \BlkEntities entities the naive comparison
space is \BlkPairSpace pairs, so blocking is not optional.

\paragraph{One identity, stated up front.} Token blocking and first-token blocking key on
the \texttt{core()} token function, and the invisible stratum is \emph{defined} as an
empty intersection under that same function. Their \BlkTokenInvisible recall on that
stratum is therefore a tautology, provable without running anything, and we claim no
credit for it. Presenting it as an empirical discovery would be circular, and a reader
who works on blocking would notice immediately.

The empirical question is what happens once the key is no longer that function. We
therefore evaluate \BlkNumSchemes schemes: the two \texttt{core()}-keyed ones for
reference, plus character q-grams, sorted-neighbourhood adjacency over lexicographic
order, phonetic Soundex keys, an attribute-keyed blocker on ZIP and city--state that
never looks at the name, and semantic nearest neighbours ($k=\BlkAnnK$) in a MiniLM
embedding space. Blocking is measured over the full recipient roster of \BlkEntities
entities, a naive space of \BlkPairSpace pairs, rather than only over entities appearing
in our own pairs file. Restricting to the latter would flatter reduction ratio and
pairs quality, because that roster has already been filtered to entities of interest.

\begin{table*}[t]
\centering
\small
\begin{tabular}{lrrrrrr}
\toprule
& & & & \multicolumn{3}{c}{Pair completeness by visibility} \\
\cmidrule(lr){5-7}
Blocking scheme & Candidates & PC (\%) & RR (\%) & Identical & Visible & Invisible \\
\midrule
Token blocking$^\dagger$ & 1,534,063 & 68.7 & 99.977 & 92.22 & 83.68 & 0.00 \\
First-token blocking$^\dagger$ & 636,042 & 52.9 & 99.990 & 78.70 & 44.80 & 0.00 \\
Q-gram blocking (q=4) & 9,022,384 & 71.4 & 99.862 & 97.83 & 78.84 & 2.79 \\
Sorted neighbourhood (w=20) & 2,170,180 & 71.2 & 99.967 & 99.95 & 73.48 & 2.05 \\
Phonetic (Soundex) & 411,063 & 61.3 & 99.994 & 97.32 & 35.96 & 0.25 \\
Attribute (ZIP / city-state) & 126,843 & 1.9 & 99.998 & 0.91 & 3.73 & 2.39 \\
Embedding ANN (MiniLM, k=20) & 1,761,025 & 72.5 & 99.973 & 99.89 & 78.67 & 2.89 \\
\midrule
Union: string-keyed & 11,875,526 & 75.2 & 99.818 & 100.00 & 89.98 & 3.82 \\
Union: all schemes & 13,074,120 & 76.2 & 99.800 & 100.00 & 91.03 & 6.77 \\
\bottomrule
\end{tabular}
\caption{Blocking over the full recipient roster, not merely the entities appearing in
our pairs file. PC is pair completeness, the recall of true links into the candidate
set; RR is reduction ratio. Rows marked $^\dagger$ key on the same \texttt{core()}
token function that \emph{defines} the invisible stratum, so their zero recall there is
an identity and not a measurement. The informative rows are the others: character
q-grams, sorted-neighbourhood adjacency, phonetic keys, attribute keys that ignore the
name entirely, and semantic nearest neighbours in an embedding space. None of them
exceeds three percent on the invisible stratum, and unioning every scheme still leaves
the overwhelming majority of those links outside the candidate set. The failure is
therefore a property of candidate generation rather than of any particular key.}
\label{tab:blocking}
\end{table*}


Table~\ref{tab:blocking} is, we think, the most consequential result in this paper, and
the informative rows are the ones not marked with a dagger. Relaxing the key does not
help. Character q-grams reach \BlkQgramInvisible on the invisible stratum, sorted
neighbourhood \BlkSnInvisible, and phonetic keys \BlkPhonInvisible. Abandoning the name
altogether and keying on address gives \BlkAttrInvisible. Semantic nearest neighbours
reach \BlkAnnInvisible, and they are the strongest single scheme overall at \BlkAnnPC
pair completeness and emphatically not a function of token overlap. Unioning every scheme, which
bounds the whole family, gives \BlkUnionPC overall and \BlkUnionInvisible on the
invisible stratum.

So the barrier is not the choice of key, and it is not that practitioners have simply
picked the wrong one. \BlkInvisibleLost of invisible links are absent from the candidate
set produced by every blocking strategy we could construct, including semantic and
non-name ones. Recent work has begun to measure pair-completeness disparity across
subgroups as a first-class property of blockers~\citep{gagliardelli2024blockingbias};
what we add is a task where the disparity is near-total and where the affected subgroup
is the one that carries the commercial value.

This reframes the matching results that follow. It is tempting to read a poor
invisible-stratum $F_1$ as a call for a better matcher: a stronger encoder, a larger
language model, better fine-tuning. That response cannot work, because in a deployed
pipeline the pair is never presented to the matcher at all. Progress requires
intervening at candidate generation, with evidence retrieved from outside the records:
ownership filings, news, or a knowledge base. Reranking cannot recover a candidate that
was never generated.

\section{Matching Baselines}
\label{sec:results}

\subsection{Protocol}

We evaluate \NumStringMethods string-similarity methods, all dependency-free so the
benchmark reproduces from a bare Python: exact match on normalised strings, Jaccard
overlap of distinctive tokens, a character sequence ratio standing in for edit distance,
and character 3-gram TF--IDF cosine with idf fitted on the training split only. We add a
neural baseline: cosine similarity between \EmbModel sentence
embeddings~\citep{zeakis2023pretrained}, encoding \EmbNames distinct names into \EmbDim dimensions.
For every method the decision threshold is chosen to maximise $F_1$ on validation and
applied unchanged to test; selecting it on test would inflate every figure here.

\subsection{The aggregate hides the failure}

Table~\ref{tab:baselines} is the central matching result. Read the overall column alone
and the task looks close to solved: \BestMethod reaches \BestOverallF $F_1$, with
precision \BestOverallPrec and recall \BestOverallRecall. Read across the strata and a
different picture emerges. On identical pairs it scores \BestIdenticalF, unsurprising
because the pairs are string-equal. On visible pairs, \BestVisibleF. On invisible pairs,
\BestInvisibleF, and no method exceeds \AnyInvisibleFMax.

\begin{table*}[t]
\centering
\small
\begin{tabular}{lrrrrr}
\toprule
& & & \multicolumn{3}{c}{Recall (\%) by name visibility} \\
\cmidrule(lr){4-6}
Method & AP & Prec. & Identical & Visible & Invisible \\
\midrule
Exact (normalised) & 65.4 & 97.3 & 100.0 & 0.0 & 0.0 \\
Token Jaccard & 73.5 & 91.8 & 99.2 & 43.5 & 0.0 \\
Sequence ratio & 74.7 & 94.7 & 100.0 & 15.8 & 3.8 \\
TF--IDF char 3-gram & 78.7 & 89.6 & 100.0 & 49.9 & 4.2 \\
\midrule
MiniLM embedding & -- & 88.6 & 100.0 & 64.3 & 4.7 \\
\midrule
\emph{Answer yes to everything} & -- & 25.0 & 100.0 & 100.0 & 100.0 \\
\bottomrule
\end{tabular}
\caption{Baselines on the \textsc{CorpFam} test split, with thresholds selected on
validation and applied unchanged to test. \textbf{Recall is reported per stratum
rather than $F_1$}, because the strata have positive rates of roughly 97\%, 10\% and
17\% respectively, so a per-stratum $F_1$ largely reflects that base rate rather than
the method: the final row shows a classifier that answers yes to every pair, which
attains perfect recall everywhere and an $F_1$ ranging from 18 to 99 across strata.
Recall depends only on a stratum's positives and is therefore comparable across them.
Average precision is threshold-free and computed over the whole split, since precision
depends on the shared negative pool. Every method recovers essentially all identical
pairs and almost no invisible ones.}
\label{tab:baselines}
\end{table*}


The gap between the identical and invisible strata is \VisibilityGap $F_1$ points, and a
single aggregate conceals all of it. Exact match and token Jaccard score \emph{zero} on
the invisible stratum, which is not a deficiency of those methods but a definitional
consequence: the stratum is constructed so no shared token exists, and methods operating
on shared tokens have nothing to operate on.

The sentence encoder is the strongest method overall and is meaningfully better on the
visible stratum, which is what one expects: it handles paraphrase and abbreviation that
token overlap misses. On the invisible stratum it gains \EmbOverStringInvisible points
over the best string method. It recovers essentially none of the gap.

\subsection{False positives concentrate where deployment lives}

Table~\ref{tab:negatives} decomposes false positives by negative grade. Random negatives
are nearly free for every method. False positives concentrate in the hard-string and
same-block grades, confirming that the difficulty sits exactly where a deployed pipeline
encounters it, after blocking and among candidates that already look alike, and that a
benchmark using random negatives alone would report a substantially rosier precision than
any practitioner will observe.

\begin{table}[t]
\centering
\small
\begin{tabular}{lrrrrr}
\toprule
Negative kind & Exact & Jaccard & SeqRatio & TF--IDF & MiniLM \\
\midrule
Random & 0.0 & 0.0 & 0.0 & 0.0 & 0.0 \\
Hard string & 1.0 & 3.6 & 1.6 & 4.6 & 6.5 \\
Same block & 0.5 & 2.2 & 1.7 & 3.2 & 2.6 \\
\bottomrule
\end{tabular}
\caption{False-positive rate (\%) by negative kind on the test split. Random
negatives are nearly free; the difficulty is concentrated in negatives that share a
blocking key or distinctive tokens, which is what a deployed pipeline actually sees
after blocking.}
\label{tab:negatives}
\end{table}


\subsection{Spend weighting}

Procurement value is heavily concentrated: across all \BenchPositives links the top 1\%
carry \SpendTopOnePct of obligated value and the top 10\% carry \SpendTopTenPct. (An
earlier version of this figure was computed on the per-entity API sample, which an
interrupted sorted fetch had left almost entirely composed of identifiers beginning with
a single letter; it is not a population sample and is superseded here by the full link
set.) A method that resolves many small suppliers while missing a few large ones is not
equivalent to one that does the reverse, though $F_1$ scores them identically. We
therefore also report recall weighted by the obligated amount on the child entity; for
\BestMethod, spend-weighted recall on test is \BestSpendRecall against an unweighted
recall of \BestOverallRecall.

\subsection{Do attributes rescue the invisible stratum?}
\label{sec:attributes}

The sharpest objection to everything above is circularity: pairs are labelled invisible
precisely because their names share nothing, and then name-based methods are shown to fail
on them. Real entity resolution uses addresses, and the established benchmarks are
multi-attribute for that reason. The objection deserves a measurement.

The award archive carries address, city, state, ZIP, country, phone and a
doing-business-as name for every recipient, giving attributes on \AttrEntities entities.
We train logistic regression on the same validation-selected-threshold protocol over
three feature sets: names only, attributes only, and both.

The evaluation is restricted to the \AttrSubsetTest test pairs carrying attributes on
\emph{both} sides. That restriction is necessary rather than convenient: over the full
split an attribute-only model has an all-zero feature vector for most pairs and
degenerates to predicting the majority class, producing a base-rate $F_1$ that measures
nothing. Even on the restricted subset the base rates differ by stratum, so
Table~\ref{tab:multiattribute} reports each cell against its majority-class floor.

\begin{table*}[t]
\centering
\small
\begin{tabular}{lrrrrrr}
\toprule
& \multicolumn{3}{c}{Overall} & \multicolumn{3}{c}{$F_1$ by name visibility} \\
\cmidrule(lr){2-4}\cmidrule(lr){5-7}
Feature set & P & R & $F_1$ & Identical & Visible & Invisible \\
\midrule
Name features only & 79.1 & 47.6 & 59.5 & 96.9 & 51.3 & 31.2 \\
Attributes only & 40.1 & 98.4 & 57.0 & 96.2 & 42.6 & 59.9 \\
Name $+$ attributes & 49.5 & 81.7 & 61.7 & 96.9 & 49.0 & 58.4 \\
\bottomrule
\end{tabular}
\caption{Supervised logistic regression over name and attribute features, evaluated on
the subset of test pairs carrying attributes (address, city, state, ZIP, country,
phone, doing-business-as name) on \emph{both} sides. Attribute evidence does supply
something names cannot: on the invisible stratum, where name features score zero by
construction, attributes reach the high twenties. That is real signal and it is worth
having, but it leaves the stratum unsolved. The deeper obstacle is coverage: attributes
exist on both sides for only a small fraction of pairs, because a parent is frequently a
holding entity that never transacts and therefore has no address at all.}
\label{tab:multiattribute}
\end{table*}


The honest answer is that attributes do not rescue the stratum, and the reason is
structural rather than a matter of feature engineering.

\textbf{The parent usually has no attributes to compare.} Of \ParentsTotal distinct
parents, only \ParentsTransacting (\ParentsTransactingPct) ever appear as an award
recipient in their own right. For the remaining \ParentsNoFootprintPct there is no
address, no phone and no location anywhere in the source, because a holding company that
never transacts has no operating footprint to record. This is a property of what a parent
\emph{is}, not a gap that better sourcing would close, and it caps attribute-based
matching at a small minority of pairs before any model is trained.

\textbf{Where attributes do exist, they mostly disagree.} Only
\AttrInvisibleSameAddrPct of true invisible pairs with addresses on both sides actually
share one. Subsidiaries operate from their own premises. That is largely what makes
them subsidiaries rather than branches.

Consequently the attribute-only and combined models on the restricted subset do not
separate from their majority-class baselines by a margin the sample size can support
($n=\AttrSubsetTest$, of which a minority are invisible positives). We report the
comparison rather than a headline number, and we note the one robust qualitative finding:
name features, which carry the identical and visible strata, contribute nothing on the
invisible stratum and can be actively misleading there.

\section{Family-Level Clustering}
\label{sec:clustering}

Pairwise linkage is not what a procurement or credit system needs. The deliverable is a
partition: every supplier record assigned to a family so spend can be rolled up. Pairwise
$F_1$ does not measure that, because pairwise errors compound through transitivity: one
false link merges two families and corrupts both rollups.

We therefore evaluate the end-to-end job over the \ClusFamilies families with at least two
children, scoring connected components of the thresholded similarity graph with pairwise
metrics, B-cubed metrics, and exact family recovery. Two conditions are reported:
\emph{blocked}, where candidates come from token blocking as in deployment, and
\emph{oracle}, where every within-split pair is scored.

\begin{table}[t]
\centering
\footnotesize
\setlength{\tabcolsep}{4pt}
\begin{tabular}{lrrrrrr}
\toprule
& & & \multicolumn{3}{c}{B-cubed} & Exact \\
\cmidrule(lr){4-6}
Condition & Cands & Pair $F_1$ & P & R & $F_1$ & fam.\ (\%) \\
\midrule
Blocked & 5,080 & 32.2 & 85.8 & 45.4 & 59.3 & 16.6 \\
Oracle & 536,130 & 39.4 & 84.3 & 45.1 & 58.8 & 14.9 \\
\bottomrule
\end{tabular}
\caption{Family-level clustering over the 902 families with
two or more children, scored on test with thresholds chosen on validation.
\emph{Blocked} takes candidates from token blocking as in deployment;
\emph{Oracle} scores every within-split pair. Removing blocking raises pairwise
$F_1$ but \emph{lowers} exact family recovery, because the additional candidates
introduce false links that merge distinct families through transitivity. Only a small
minority of families are recovered exactly under either condition, and exact recovery
is what a spend rollup depends on.}
\label{tab:clustering}
\end{table}


Under blocking the system reaches \ClusBlockedBcubed B-cubed $F_1$ and recovers
\ClusBlockedExact of families exactly. Removing blocking raises pairwise $F_1$ from
\ClusBlockedPairF to \ClusOraclePairF, yet exact family recovery \emph{falls} to
\ClusOracleExact and B-cubed $F_1$ to \ClusOracleBcubed. The extra candidates introduce
false links, and transitive closure propagates each into a merger of two distinct
families.

We do not want to overstate this. The comparison rests on \ClusTestFamilies test
families, the confidence intervals on exact recovery overlap, and the direction of the
pairwise and partition metrics diverges rather than one dominating. The defensible claim
is the narrow one, and it is still worth making: pairwise $F_1$ and exact partition
recovery are not monotonically related, so a pipeline tuned on pairwise quality can move
the deliverable backwards. Any evaluation of this task that reports only pairwise
metrics is measuring something other than what the task exists to produce.

\section{Discussion}
\label{sec:discussion}

\paragraph{The invisible stratum measures knowledge, not similarity.} The stratum is
defined by the absence of shared distinctive tokens. For a pair such as (\textsc{Blue
Aerospace LLC}, \textsc{Heico Corp}) there is no function of the two strings that returns
the correct answer, because the strings do not contain it. The relationship exists in an
SEC filing, as Section~\ref{sec:edgar} shows by going and reading one. Whatever supplies
the answer has to come from outside the pair, which reframes corporate-family resolution
as a retrieval and evidence-aggregation problem rather than a matching problem. That
connects it to work on grounded factual
verification~\citep{wei2024simpleqa,wei2025browsecomp,chen2025browsecompplus} rather than
to the entity-matching literature it is usually filed under. We did not test whether a
pretrained model already holds these facts, and we make no claim either way; that is a
separate experiment with its own contamination problem, since the filings are on the open
web.

\paragraph{The intervention point is candidate generation, not ranking.}
Section~\ref{sec:blocking} sharpens this considerably. Because \BlkInvisibleLost of
invisible links never enter the candidate set, work at the matching stage is spent on
pairs that a deployed system would never adjudicate. A retrieval-based blocker, one
that consults ownership filings to propose candidates that share no string evidence,
would change the achievable ceiling in a way that no matcher can.

\paragraph{What this implies for evaluation practice.} The specific failure documented
here, a defensible aggregate concealing near-zero performance on the subset that motivates
the task, is not special to corporate families. It arises whenever a benchmark's positives
are heterogeneous in difficulty and the easy majority correlates with a surface feature.
We would suggest benchmark builders report a stratification along whatever axis makes the
easy cases easy, and that reviewers ask for it. Our own negative-sampling error in
Section~\ref{sec:splits} is a second, related caution: when the two sides of a relation are
different kinds of object, a model can learn the object type instead of the relation, and
an aggregate metric will not reveal it.

\section{Self-Assessment of Relevance}

\textsc{CorpFam} is a data-management contribution in three respects, and we state them
explicitly because a benchmark drawn from procurement data could be mistaken for an
applied data-cleaning exercise.

First, the object of study is entity resolution, a core data-integration problem with a
long history at this venue, and the artifact is a reusable benchmark with held-out splits
of the kind the Experiment, Analysis \& Benchmark category exists to publish.

Second, the central finding is about \emph{blocking}, the indexing and candidate
generation step that is a database problem rather than a modelling one. We show that the
standard family of blocking keys has a systematic blind spot that aggregate pair
completeness hides, and we quantify it per stratum. That is a statement about the
architecture of resolution pipelines, not about any particular classifier.

Third, the paper contributes to evaluation methodology: a stratification that makes
otherwise incomparable positives comparable, a spend-weighted metric reflecting that value
is concentrated, a family-level clustering task showing pairwise and partition quality can
move in opposite directions, and a quantified ceiling on the ground truth itself.

\section{Limitations}

Our ground truth is a single jurisdiction and a single fiscal year, so families are those
visible in US federal contracting in FY2025; multinational structures appear only through
their US-registered entities. The parent relation as recorded is one hop, so we evaluate
child-to-ultimate-parent linkage, not intermediate holding chains. Label noise is bounded
but real, as Section~\ref{sec:quality} quantifies, and our hand adjudication covers large
families exhaustively and small families not at all. The Exhibit~21 corroboration covers
SEC registrants only. The multi-attribute baseline is measured on the subset with
attributes on both sides, which is a minority of pairs. Finally, our baselines establish
the shape of the problem and a floor rather than the state of the art; we release the
benchmark so stronger methods can be measured against the stratification rather than
against the aggregate.

\section{Conclusion}

Corporate-family resolution is treated as a variant of entity matching and evaluated with
a single aggregate metric. We built a public benchmark of \BenchPairs pairs over
\BenchFamilies families from self-reported registry data and showed that this practice is
misleading in two distinct ways. At the matching stage an overall $F_1$ of \BestOverallF
coexists with at most \AnyInvisibleFMax on pairs where the relationship is not visible in
the names. More fundamentally, the standard pipeline never gets that far: no string-keyed
blocking scheme proposes those candidates, so \BlkInvisibleLost of them are lost before
matching begins. The links are real: \EdgarInvisibleConfirmPct of them are corroborated
by the parent's own SEC filings. Conventional attribute evidence mitigates the
difficulty without resolving it. The pairs that matter commercially are the ones current pipelines
cannot represent at all, and recognising that requires reporting the stratum rather than
the average.

\section*{Data and Code Availability}

The benchmark, the adjudicated exclusion list with per-family reasons, and all code
required to rebuild the dataset and reproduce every number in this paper are released. All
source data is public United States federal award data~\citep{usaspending_api} and SEC
EDGAR filings. No confidential or proprietary information was used. Every numeric claim in
this manuscript is generated from a results file by \texttt{make\_tables.py} and verified
by \texttt{check\_manuscript.py}; no figure is transcribed by hand. The build is seeded
(\BenchSeed) and deterministic.
